\documentclass[aspectratio=169,11pt]{beamer}

% Shared style, ported from the UF Animal Sciences 2026 talk (ib-tutorials.sty at the repo root).
\makeatletter\def\input@path{{./}{../}}\makeatother
\usepackage{ib-tutorials}

\usepackage[utf8]{inputenc}
\usepackage{amsmath}
\usepackage{amssymb}
\usepackage{graphicx}
% Resolve ../pictures/... whether compiling in this folder or from the repo root.
\graphicspath{{./}{../}{2-Solid-Mechanics-Under-Construction/}}

\title{Mathematical Description of Shape Changes in Solids}
\author{Laura A. Miller}
\institute{Department of Mathematics, University of Arizona}
\date{Fluid Dynamics and IBM Theory \textbullet\ Module 2}

\begin{document}

\begin{frame}[plain,b]
\maketutorialtitle
\end{frame}

\begin{frame}{Introduction}
  \begin{itemize}
    \item   The purpose of this chapter is to summarize the equations that govern the response of solids to mechanical or thermal loading.
    \item   We will focus on the mathematical description of shape changes in a solid.
    \item   This includes displacement and velocity fields, and deformation gradient tensors.
  \end{itemize}
  \vspace{3mm}
  {\scriptsize\color{DarkGray}This deck closely follows A. F. Bower, \emph{Applied Mechanics of Solids} (CRC Press, 2010), \S2.1 (\url{https://solidmechanics.org}); text, equations and figures are adapted from that section, which is made freely available for non-commercial educational use. Notation: $\mathbf{x}$ = reference position, $\mathbf{y}$ = deformed position (Bower's convention); deck 2-2 writes $\mathbf{w}$ and decks 4-x write $\boldsymbol{\chi}(\mathbf{X},t)$ for the same map.\par}
\end{frame}

\begin{frame}{Mathematical description of shape changes in solids}
    \begin{itemize}
        \item   In this section, we list the various mathematical formulas that are used to characterize shape changes in solids (and in fluids).
        \item   The formulas might look scary at first, but they are mostly just definitions.
        \item   Shape changes near a point can always be characterized by six numbers.
        \item   (Why six: the deformation gradient $\mathbf{F}$ has nine components, but three of them describe a local rigid rotation, which does not change shape. A symmetric strain tensor has the remaining six.)
    \end{itemize}
\end{frame}

\begin{frame}{Mathematical description of shape changes in solids (continued)}
    \begin{itemize}
        \item   These could be the six independent components of the Lagrangian strain, Eulerian strain, the left or right stretch tensors, or your own favorite deformation measure (these are developed in deck 4-2, IBFE Solids Prep).
        \item   Given the complete set of six numbers for any one deformation measure, you can always calculate the components of other strain measures.
        \item   The reason that so many different deformation measures exist is partly that different material models adopt different strain measures, and partly because each measure is useful for describing a particular type of shape change.
    \end{itemize}
\end{frame}

\begin{frame}{The Displacement and Velocity Fields}
\begin{columns}[c,onlytextwidth]
\begin{column}{0.52\textwidth}
  \begin{itemize}
    \item   The displacement vector $\mathbf{u}(\mathbf{x}, t)$ describes the motion of each point in the solid.
    \item   The displacement vector is defined as:
    $$ \mathbf{y} = \mathbf{x} + \mathbf{u}(\mathbf{x}, t) $$
    i.e.\ $\mathbf{u} = \mathbf{y} - \mathbf{x}$, where $\mathbf{x}$ is the position of a material particle in the undeformed (reference) solid and $\mathbf{y}$ its position at time $t$.
    \item   In index notation:
    $$ y_{i}=x_{i}+u_{i}(x_{1},x_{2},x_{3},t) $$
    \item   The displacement field completely specifies the change in shape of the solid.
  \end{itemize}
\end{column}
\begin{column}{0.46\textwidth}
\begin{center}
   \includegraphics[width=\linewidth]{../pictures/2-3/image004.png}
\end{center}
{\centering\tiny\color{DarkGray}Figure: A. F. Bower, \emph{Applied Mechanics of Solids} (CRC Press, 2010), \S2.1.1, solidmechanics.org (free for non-commercial educational use).\par}
\end{column}
\end{columns}

\end{frame}

\begin{frame}{Velocity Field}
  \begin{itemize}
    \item   The velocity field describes the motion of the solid.
    \item   The velocity field is defined as the time derivative of the displacement field:
    $$ v_{i}(x_{k},t)=\frac{\partial y_{i}}{\partial t}=\frac{\partial u_{i}(x_{k},t)}{\partial t} \qquad (\mathbf{x} = \text{const}) $$
    \item   ``$\mathbf{x}$ = const'' means the derivative is taken with the \emph{reference} position $\mathbf{x}$ held fixed, i.e.\ following one material particle. This is the \textbf{material (Lagrangian) velocity}; the second equality holds because $\mathbf{x}$ does not depend on $t$.
    \item   Similarly the acceleration is $a_i = \partial^2 u_i/\partial t^2$ at fixed $\mathbf{x}$, the quantity that appears in the equations of motion (deck 2-1).
    \item   In a fluid one usually works with the spatial (Eulerian) velocity field $\mathbf{v}(\mathbf{y},t)$; the two descriptions are linked by $\mathbf{v}_{\rm material}(\mathbf{x},t) = \mathbf{v}_{\rm spatial}(\mathbf{y}(\mathbf{x},t),t)$ (decks 0-1 and 4-2).
  \end{itemize}
\end{frame}

\begin{frame}{Examples of Simple Deformations}
\begin{columns}[c,onlytextwidth]
\begin{column}{0.46\textwidth}
  \begin{itemize}
    \item   \textbf{Volume preserving uniaxial extension}
    $$
    \begin{aligned}
      y_{1} &= \lambda x_{1} \\
      y_{2} &= x_{2} / \sqrt{\lambda} \\
      y_{3} &= x_{3} / \sqrt{\lambda}
    \end{aligned}
    $$

  \end{itemize}
\end{column}
\begin{column}{0.52\textwidth}
\begin{center}
   \includegraphics[width=0.9\linewidth]{../pictures/2-3/image008.png}
\end{center}
{\centering\tiny\color{DarkGray}Figure: Bower (2010), \S2.1.1, solidmechanics.org (free for non-commercial educational use).\par}
\end{column}
\end{columns}
\end{frame}

\begin{frame}{Examples of Simple Deformations}
\begin{columns}[c,onlytextwidth]
\begin{column}{0.46\textwidth}
  \begin{itemize}
    \item   \textbf{Simple shear}
    $$
    \begin{aligned}
      y_{1} &= x_{1} + \tan(\theta) x_{2} \\
      y_{2} &= x_{2} \\
      y_{3} &= x_{3}
    \end{aligned}
    $$
  \end{itemize}
\end{column}
\begin{column}{0.52\textwidth}
\begin{center}
   \includegraphics[width=0.9\linewidth]{../pictures/2-3/image011.png}
\end{center}
{\centering\tiny\color{DarkGray}Figure: Bower (2010), \S2.1.1, solidmechanics.org (free for non-commercial educational use).\par}
\end{column}
\end{columns}
\end{frame}

\begin{frame}{Examples of Simple Deformations (continued)}
\begin{columns}[c,onlytextwidth]
\begin{column}{0.46\textwidth}
  \begin{itemize}
    \item   \textbf{Rigid rotation through angle \(\theta\) about \(e_3\) axis}
    $$
    \begin{aligned}
      y_{1} &= x_{1} \cos(\theta) - x_{2} \sin(\theta) \\
      y_{2} &= x_{2} \cos(\theta) + x_{1} \sin(\theta) \\
      y_{3} &= x_{3}
    \end{aligned}
    $$
  \end{itemize}
\end{column}
\begin{column}{0.52\textwidth}
  \begin{center}
   \includegraphics[width=0.9\linewidth]{../pictures/2-3/rigid-rot.png}
\end{center}
{\centering\tiny\color{DarkGray}Figure: Bower (2010), \S2.1.1, solidmechanics.org (free for non-commercial educational use).\par}
\end{column}
\end{columns}
\end{frame}

\begin{frame}{Examples of Simple Deformations (continued)}
\begin{columns}[c,onlytextwidth]
\begin{column}{0.55\textwidth}
  \begin{itemize}
    \item   \textbf{General rigid rotation about the origin}
    $$ \mathbf{y} = \mathbf{R} \cdot \mathbf{x} \quad \text{or} \quad y_{i}=R_{ij}x_{j} $$
    where \(\mathbf{R}\) must satisfy \(\mathbf{R} \cdot \mathbf{R}^T = \mathbf{R}^T \cdot \mathbf{R} = \mathbf{I}\) and \(\det(\mathbf{R}) = +1\) (\(\mathbf{R}\) is proper orthogonal). \(\mathbf{I}\) is the identity tensor
    with components
    $$
    \delta_{ik}=\begin{cases}
      1, & i=k \\
      0, & i \ne k
    \end{cases}
    $$
  \end{itemize}
\end{column}
\begin{column}{0.43\textwidth}
  \begin{center}
   \includegraphics[width=\linewidth]{../pictures/2-3/gen-rigid-rot.png}
\end{center}
{\centering\tiny\color{DarkGray}Figure: Bower (2010), \S2.1.1, solidmechanics.org (free for non-commercial educational use).\par}
\end{column}
\end{columns}
\end{frame}

\begin{frame}{Examples of Simple Deformations (continued)}
  \begin{itemize}
    \item   \textbf{General rigid rotation about the origin}
    Alternatively, a rigid rotation through angle \(\theta\) (with right hand screw convention) about an axis through the origin that is parallel to a unit vector \(\mathbf{n}\) can be written as
    $$ \mathbf{y} = \cos\theta \mathbf{x} + (1 - \cos\theta)(\mathbf{n} \cdot \mathbf{x})\mathbf{n} + \sin\theta (\mathbf{n} \times \mathbf{x}) $$
    \item   The components of R are thus
    $$ R_{ij} = \cos\theta \delta_{ij} + (1 - \cos\theta)n_i n_j + \sin\theta \epsilon_{ikj} n_k \;\big(= \cos\theta \delta_{ij} + (1 - \cos\theta)n_i n_j - \sin\theta \epsilon_{ijk} n_k\big)$$
    where \(\epsilon_{ijk}\) is the permutation symbol, satisfying
    $$
    \epsilon_{ijk}=\begin{cases}
      1 & i,j,k = 1,2,3; \quad 2,3,1 \text{ or } 3,1,2 \\
      -1 & i,j,k = 3,2,1; \quad 2,1,3 \text{ or } 1,3,2 \\
      0 & \text{otherwise}
    \end{cases}
    $$
    \item   Check: with $\mathbf{n}=\mathbf{e}_3$ this reduces to the rotation about $\mathbf{e}_3$ on the previous slide ($R_{12} = \sin\theta\,\epsilon_{132} = -\sin\theta$, $R_{21} = \sin\theta\,\epsilon_{231} = +\sin\theta$).
  \end{itemize}
\end{frame}

\begin{frame}{Examples of Simple Deformations (continued)}
  \begin{itemize}
    \item   \textbf{General homogeneous deformation}
    $$
    \begin{aligned}
      y_1 &= A_{11} x_1 + A_{12} x_2 + A_{13} x_3 + c_1 \\
      y_2 &= A_{21} x_1 + A_{22} x_2 + A_{23} x_3 + c_2 \\
      y_3 &= A_{31} x_1 + A_{32} x_2 + A_{33} x_3 + c_3
    \end{aligned}
    $$
    or
    $$ \mathbf{y} = \mathbf{A} \cdot \mathbf{x} + \mathbf{c} \quad y_i = A_{ij} x_j + c_i $$
    where \(A_{ij}\) are constants.
    \item   The physical significance of a homogeneous deformation is that all straight lines in the solid remain straight under the deformation.
    \item   Thus, every point in the solid experiences the same shape change. All the deformations listed above are examples of homogeneous deformations.
    \item   ($\mathbf{c}$ is a rigid translation; the shape change is carried entirely by the constant matrix $\mathbf{A}$, which is the deformation gradient $\mathbf{F}$ defined below.)
  \end{itemize}
\end{frame}

\begin{frame}{Worked Examples: $\mathbf{F}$ and $J$ for the Simple Deformations}
  \begin{itemize}
    \item Using $F_{ij} = \partial y_i/\partial x_j$ (defined on the next slides) and $J=\det\mathbf{F}$ (end of this deck):
  \end{itemize}
  \[
  \renewcommand{\arraystretch}{1.15}
  \begin{array}{lcc}
  \text{deformation} & \mathbf{F} & J=\det\mathbf{F}\\ \hline
  \text{uniaxial extension} & \mathrm{diag}\big(\lambda,\ \lambda^{-1/2},\ \lambda^{-1/2}\big) & \lambda\cdot\lambda^{-1/2}\cdot\lambda^{-1/2} = 1\\[3pt]
  \text{simple shear} & \begin{bmatrix}1&\tan\theta&0\\0&1&0\\0&0&1\end{bmatrix} & 1\\[14pt]
  \text{rotation about } \mathbf{e}_3 & \begin{bmatrix}\cos\theta&-\sin\theta&0\\ \sin\theta&\cos\theta&0\\0&0&1\end{bmatrix} & \cos^2\theta+\sin^2\theta = 1\\[14pt]
  \text{homogeneous} & \mathbf{A} & \det\mathbf{A}
  \end{array}
  \]
  \begin{itemize}
    \item All three named examples preserve volume ($J=1$), but only the rotation preserves shape: $\mathbf{F}^T\mathbf{F}=\mathbf{I}$ for the rotation, whereas $\mathbf{F}^T\mathbf{F}\neq\mathbf{I}$ for extension and shear.
  \end{itemize}
\end{frame}

\begin{frame}{Displacement Gradient Tensor}
    \begin{itemize}
        \item The displacement gradient tensor, denoted as \(\mathbf{u} \otimes \nabla\), is a second-order tensor.
        \item It is represented in component form as:
            \[
            \begin{bmatrix}
                \frac{\partial u_1}{\partial x_1} & \frac{\partial u_1}{\partial x_2} & \frac{\partial u_1}{\partial x_3} \\
                \frac{\partial u_2}{\partial x_1} & \frac{\partial u_2}{\partial x_2} & \frac{\partial u_2}{\partial x_3} \\
                \frac{\partial u_3}{\partial x_1} & \frac{\partial u_3}{\partial x_2} & \frac{\partial u_3}{\partial x_3}
            \end{bmatrix}
            \]
            where \(u_1\), \(u_2\), and \(u_3\) are the components of the displacement vector \(\mathbf{u}\), and \(x_1\), \(x_2\), and \(x_3\) are the components of the position vector \(\mathbf{x}\).
        \item Bower writes this as \(\mathbf{u} \otimes \nabla\); deck 2-2 writes \(\nabla \mathbf{u}\) with \((\nabla \mathbf{u})_{ij} = \partial u_i / \partial x_j\); they are the same object.
        \item Caution: a few texts define \(\nabla\mathbf{u}\) as the \emph{transpose}, \((\nabla\mathbf{u})_{ij} = \partial u_j/\partial x_i\) (sometimes written \(\nabla\otimes\mathbf{u}\)). Always check which index is differentiated; these decks use \(\partial u_i/\partial x_j\) throughout.
    \end{itemize}
\end{frame}

\begin{frame}{Deformation Gradient Tensor}
    \begin{itemize}
        \item The deformation gradient tensor, denoted as \(\mathbf{F}\), is a second-order tensor.
        \item It relates infinitesimal line elements in the undeformed and deformed configurations.
        \item It is defined as:
            \[
            \mathbf{F} = \mathbf{I} + \mathbf{u} \otimes \nabla
            \]
            where \(\mathbf{I}\) is the identity tensor, \(\mathbf{u}\) is the displacement vector, and \(\nabla\) is the gradient operator.
            \end{itemize}
            \end{frame}

\begin{frame}{Deformation Gradient Tensor}
\begin{itemize}
        \item In component form, this tensor is represented as:
            \[
            F_{ij} = \delta_{ij} + \frac{\partial u_i}{\partial x_j} = \frac{\partial y_i}{\partial x_j}
            \]
            or in matrix form:
            \[
            \begin{bmatrix}
                1 + \frac{\partial u_1}{\partial x_1} & \frac{\partial u_1}{\partial x_2} & \frac{\partial u_1}{\partial x_3} \\
                \frac{\partial u_2}{\partial x_1} & 1 + \frac{\partial u_2}{\partial x_2} & \frac{\partial u_2}{\partial x_3} \\
                \frac{\partial u_3}{\partial x_1} & \frac{\partial u_3}{\partial x_2} & 1 + \frac{\partial u_3}{\partial x_3}
            \end{bmatrix}
            =
            \begin{bmatrix}
                \frac{\partial y_1}{\partial x_1} & \frac{\partial y_1}{\partial x_2} & \frac{\partial y_1}{\partial x_3} \\
                \frac{\partial y_2}{\partial x_1} & \frac{\partial y_2}{\partial x_2} & \frac{\partial y_2}{\partial x_3} \\
                \frac{\partial y_3}{\partial x_1} & \frac{\partial y_3}{\partial x_2} & \frac{\partial y_3}{\partial x_3}
            \end{bmatrix}
            \]
            where \(u_i\) are the components of the displacement vector \(\mathbf{u}\), \(x_j\) are the components of the position vector \(\mathbf{x}\) in the undeformed configuration, \(y_i\) are the components of the position vector \(\mathbf{y}\) in the deformed configuration, and \(\delta_{ij}\) is the Kronecker delta.
    \end{itemize}
\end{frame}

\begin{frame}{Kronecker delta symbol}
    \begin{itemize}
        \item  The Kronecker delta symbol is defined as:
        $$
        \delta_{ik}=\begin{cases}
          1, & i=k \\
          0, & i \ne k
        \end{cases}
        $$
    \end{itemize}
\end{frame}

\begin{frame}{Displacement Gradient Tensor Continued}
  \begin{itemize}
    \item   Formally, the gradient of a vector field \(\mathbf{u}(\mathbf{x})\) is defined so that
    $$ [\mathbf{u} \otimes \nabla] \cdot \mathbf{n} = \lim_{\alpha \to 0} \frac{\mathbf{u}(\mathbf{x} + \alpha \mathbf{n}) - \mathbf{u}(\mathbf{x})}{\alpha} $$
    (for more details see Appendix B of https://solidmechanics.org/), but in practice the component formula \(\frac{\partial u_i}{\partial x_j}\) is more useful.
    \item   Note also that
    $$ \mathbf{y} \otimes \nabla = (\mathbf{x} + \mathbf{u}(\mathbf{x})) \otimes \nabla = \mathbf{F} $$
    or
    $$ \frac{\partial y_i}{\partial x_j} = \frac{\partial}{\partial x_j} (x_i + u_i) = \delta_{ij} + \frac{\partial u_i}{\partial x_j} = F_{ij} $$
    The rules of differentiation using index notation are described in more detail in Appendix C of https://solidmechanics.org/.
  \end{itemize}
\end{frame}

\begin{frame}{Deformation Gradient Tensor and Infinitesimal Line Elements}
\begin{columns}[c,onlytextwidth]
\begin{column}{0.50\textwidth}
  \begin{itemize}
    \item   The concepts of displacement gradient and deformation gradient are introduced to quantify the change in shape of infinitesimal line elements in a solid body.
    \item   Imagine drawing a straight line on the undeformed configuration of a solid.
    \item   The line would be mapped to a smooth curve on the deformed configuration.
  \end{itemize}
\end{column}
\begin{column}{0.48\textwidth}
    \begin{center}
   \includegraphics[width=\linewidth]{../pictures/2-3/infinitesimal.png}
\end{center}
{\centering\tiny\color{DarkGray}Figure: Bower (2010), \S2.1.2, solidmechanics.org (free for non-commercial educational use).\par}
\end{column}
\end{columns}
\end{frame}

\begin{frame}{Deformation Gradient Tensor and Infinitesimal Line Elements}
  \begin{itemize}
    \item   Focus on a line segment dx, much shorter than the radius of curvature of this curve.
    \item   The segment would be straight in the undeformed configuration, and also (almost) straight in the deformed configuration.
    \item   Infinitesimal line segments are merely stretched and rotated by a deformation.
    \item   The infinitesimal line segments dx and dy are related by
    $$ dy = F \cdot dx \quad \text{or} \quad dy_i = F_{ik} dx_k $$
  \end{itemize}
\end{frame}

\begin{frame}{Deformation Gradient Tensor as a Matrix Equation}
    \begin{itemize}
        \item   Written out as a matrix equation, we have
        $$
        \begin{bmatrix}
          dy_1 \\
          dy_2 \\
          dy_3
        \end{bmatrix} =
        \begin{bmatrix}
          1 + \frac{\partial u_1}{\partial x_1} & \frac{\partial u_1}{\partial x_2} & \frac{\partial u_1}{\partial x_3} \\
          \frac{\partial u_2}{\partial x_1} & 1 + \frac{\partial u_2}{\partial x_2} & \frac{\partial u_2}{\partial x_3} \\
          \frac{\partial u_3}{\partial x_1} & \frac{\partial u_3}{\partial x_2} & 1 + \frac{\partial u_3}{\partial x_3}
        \end{bmatrix}
        \begin{bmatrix}
          dx_1 \\
          dx_2 \\
          dx_3
        \end{bmatrix}
        $$
    \end{itemize}
\end{frame}

\begin{frame}{Derivation of  $dy = F \cdot dx$}
  \begin{itemize}
    \item   Consider an infinitesimal line element dx in a deforming solid.
    \item   When the solid is deformed, this line element is stretched and rotated to a deformed line element dy.
    \item   If we know the displacement field in the solid, we can compute \(dy = [x + dx + u(x + dx)] - [x + u(x)]\) from the position vectors of its two endpoints.
    \item   Expand \(u_{i}(x_{k}+dx_{k})\) as a Taylor series
    $$ u_{i}(x_{k}+dx_{k}) \approx u_{i}(x_{k})+\frac{\partial u_{i}}{\partial x_{k}}dx_{k} $$
    so that
    $$ dy_{i}=x_{i}+dx_{i}+u_{i}(x_{k}+dx_{k})-(x_{i}+u_{i}(x_{k})) $$
    $$ dy_{i}=dx_{i}+\frac{\partial u_{i}}{\partial x_{k}}dx_{k}=(\delta_{ik}+\frac{\partial u_{i}}{\partial x_{k}})dx_{k} $$
    \item   We identify the term in parentheses as the deformation gradient, so
    $$ dy_{i}=F_{ik}dx_{k} $$
  \end{itemize}
\end{frame}

\begin{frame}{Inverse of Deformation Gradient}
  \begin{itemize}
    \item   The inverse of the deformation gradient \(F^{-1}\) arises in many calculations.
    \item   It is defined through
    $$ dx_{i}=F_{ik}^{-1}dy_{k} $$
    \item   or alternatively
    $$ F_{ij}^{-1}=\frac{\partial x_{i}}{\partial y_{j}} $$
  \end{itemize}
\end{frame}

\begin{frame}{Deformation gradient resulting from two successive deformations (1/3)}
\begin{columns}[c,onlytextwidth]
\begin{column}{0.50\textwidth}
\begin{itemize}
    \item \textbf{Successive Deformations:} Consider two successive deformations applied to a solid.
    \item \textbf{Mapping Line Elements:}
        \begin{itemize}
            \item   The first deformation maps an infinitesimal line element $dx$ to $dy$:
                \[
                dy = F^{(1)} \cdot dx \quad \text{or} \quad dy_i = F_{ij}^{(1)} dx_j
                \]
            \item   The second deformation maps $dy$ to $dz$:
                \[
                dz = F^{(2)} \cdot dy \quad \text{or} \quad dz_i = F_{ij}^{(2)} dy_j
                \]
        \end{itemize}
\end{itemize}
\end{column}
\begin{column}{0.48\textwidth}
\begin{figure}
    \centering
    \includegraphics[width=\linewidth]{../pictures/2-3/2-deformation.png}
\end{figure}
{\centering\tiny\color{DarkGray}Figure: Bower (2010), \S2.1.3, solidmechanics.org (free for non-commercial educational use).\par}
\end{column}
\end{columns}
\end{frame}

\begin{frame}{Deformation gradient resulting from two successive deformations (2/3)}
\begin{itemize}
    \item \textbf{Deformation Gradients:}
        \begin{itemize}
            \item   $F^{(1)}$ and $F^{(2)}$ are deformation gradients:
                \[
                 F^{(1)} = y \otimes \nabla_x, \quad F^{(2)} = z \otimes \nabla_y
                \]
            \item   In component form:
                 \[
                 F_{ij}^{(1)} = \frac{\partial y_i}{\partial x_j}, \quad F_{ij}^{(2)} = \frac{\partial z_i}{\partial y_j}
                 \]
        \end{itemize}
\end{itemize}
\end{frame}

\begin{frame}{Deformation gradient resulting from two successive deformations (3/3)}
\begin{itemize}
    \item \textbf{Cumulative Deformation Gradient:}
        \begin{itemize}
            \item   The deformation gradient $F$ that maps $dx$ directly to $dz$ is:
                \[
                dz = F \cdot dx \quad \text{with} \quad F = F^{(2)} \cdot F^{(1)}
                \]
            \item   In component form:
                \[
                dz_i = F_{ij} dx_j, \quad F_{ij} = F_{ik}^{(2)} F_{kj}^{(1)}
                \]
        \end{itemize}
    \item \textbf{Chain Rule Application:}
        \begin{itemize}
            \item   The cumulative mapping can be written as $z_i(y_j(x_k))$.
            \item   Applying the chain rule:
                \[
                dz_i = \frac{\partial z_i}{\partial y_j} \frac{\partial y_j}{\partial x_k} dx_k
                \]
        \end{itemize}
\end{itemize}

\end{frame}

\begin{frame}{The Jacobian of the deformation gradient (1/4)}
\begin{itemize}
    \item \textbf{Definition:} The Jacobian of the deformation gradient is:
        \[
        J = det(F) = det\left(\delta_{ij} + \frac{\partial u_i}{\partial x_j}\right)
        \]
    \item \textbf{Volume Change:} $J$ measures the volume change due to deformation.
\end{itemize}
\begin{figure}
    \centering
    \includegraphics[scale=0.7]{../pictures/2-3/jacobian-deformation.png}
\end{figure}
\vspace{-3mm}
{\centering\tiny\color{DarkGray}Figure: Bower (2010), \S2.1.4, solidmechanics.org (free for non-commercial educational use).\par}
\end{frame}

\begin{frame}{The Jacobian of the deformation gradient (2/4)}
\begin{itemize}
        \item \textbf{Infinitesimal Volume Element:}
        \begin{itemize}
            \item   Consider an infinitesimal volume element in the \emph{reference} configuration whose three edges are the line elements $d\mathbf{x}^{(1)}$, $d\mathbf{x}^{(2)}$, and $d\mathbf{x}^{(3)}$.
            \item   Original volume: $dV_0 = d\mathbf{x}^{(1)} \cdot (d\mathbf{x}^{(2)} \times d\mathbf{x}^{(3)}) = \epsilon_{ijk}\, dx^{(1)}_i dx^{(2)}_j dx^{(3)}_k$
        \end{itemize}
    \item \textbf{Deformed Volume Element:}
        \begin{itemize}
            \item   The element is mapped to a parallelepiped whose edges are the deformed images $d\mathbf{y}^{(1)}$, $d\mathbf{y}^{(2)}$, and $d\mathbf{y}^{(3)}$.
            \item   Deformed volume: $dV = \epsilon_{ijk}\, dy^{(1)}_i dy^{(2)}_j dy^{(3)}_k$
        \end{itemize}
    \item \textbf{Relationship between deformed and undeformed elements:}
         \begin{itemize}
            \item   $dy^{(1)}_i = F_{il} dx^{(1)}_l$,  $dy^{(2)}_j = F_{jm} dx^{(2)}_m$,  $dy^{(3)}_k = F_{kn} dx^{(3)}_n$
         \end{itemize}
\end{itemize}
\end{frame}

\begin{frame}{The Jacobian of the deformation gradient (3/4)}
\begin{itemize}
    \item \textbf{Volume Transformation:}
        \begin{itemize}
            \item   $dV = \epsilon_{ijk} F_{il} F_{jm} F_{kn}\, dx^{(1)}_l dx^{(2)}_m dx^{(3)}_n$
            \item   Using the identity $\epsilon_{ijk} A_{il} A_{jm} A_{kn} = \epsilon_{lmn} det(A)$:
                \[
                dV = det(F)\, \epsilon_{lmn}\, dx^{(1)}_l dx^{(2)}_m dx^{(3)}_n = det(F)\, dV_0
                \]
        \end{itemize}
    \item \textbf{Volume Ratio:}
        \begin{itemize}
            \item   The ratio of deformed to undeformed volume:
                \[
                \frac{dV}{dV_0} = det(F) = J
                \]
        \end{itemize}
    \item \textbf{Physical Admissibility:} For any physically admissible deformation, $J > 0$.
    \item \textbf{Incompressibility:} For incompressible materials, $J = 1$.
    \item \textbf{Density Transformation:} Mass density transforms as $\rho = \frac{\rho_0}{J}$.
\end{itemize}
\end{frame}

\begin{frame}{The Jacobian of the deformation gradient (3/4, continued)}
\begin{itemize}
    \item \textbf{Why $J>0$:} $J=0$ would crush a finite volume element to zero volume, and $J<0$ would turn it inside out (a reflection), which no continuous motion starting from $\mathbf{F}=\mathbf{I}$ (where $J=1$) can reach without passing through $J=0$.
    \item \textbf{Why $\rho=\rho_0/J$:} conservation of mass for the element, $\rho\,dV = \rho_0\,dV_0$, combined with $dV = J\,dV_0$.
    \item \textbf{Identity used above:} $\epsilon_{ijk}A_{il}A_{jm}A_{kn}$ is antisymmetric in $(l,m,n)$, so it equals $\epsilon_{lmn}$ times its value for $(l,m,n)=(1,2,3)$, which is $\epsilon_{ijk}A_{i1}A_{j2}A_{k3}=\det\mathbf{A}$.
    \item \textbf{Incompressible IBFE models} enforce $J=1$ only approximately (through a volumetric energy penalising $J\neq 1$); monitoring $J$ is a useful accuracy check (deck 4-2).
\end{itemize}
\end{frame}

\begin{frame}{The Jacobian of the deformation gradient (4/4)}
\begin{itemize}
    \item \textbf{Jacobian Derivative:} Derivative of J with respect to F:
        \[
        \frac{\partial J}{\partial F_{ij}} = J F_{ji}^{-1}
        \]
    \item In tensor form, $\dfrac{\partial J}{\partial \mathbf{F}} = J\,\mathbf{F}^{-T}$.
    \item \textbf{Where it comes from (Jacobi's formula):} $dJ = J\,\mathrm{tr}\!\left(\mathbf{F}^{-1}d\mathbf{F}\right) = J\,F^{-1}_{ji}\,dF_{ij}$. Equivalently, expand $\det\mathbf{F}$ along row $i$: $\partial J/\partial F_{ij}$ is the cofactor of $F_{ij}$, and the cofactor matrix is $J\mathbf{F}^{-T}$.
    \item \textbf{Check in 1D:} $J=F_{11}$ gives $\partial J/\partial F_{11} = 1 = J F_{11}^{-1}$.
    \item This identity is used to differentiate strain-energy functions $W(J)$ when computing the first Piola--Kirchhoff stress $\mathbf{P}=\partial W/\partial\mathbf{F}$ (deck 4-2).
\end{itemize}
\end{frame}

\begin{frame}{References}
\scriptsize
\begin{itemize}
  \item A. F. Bower, \emph{Applied Mechanics of Solids}, CRC Press, Boca Raton, 2010. doi:10.1201/9781439802489. Online: \url{https://solidmechanics.org}. \S2.1.1--2.1.4 (text, equations and all figures in this deck; free for non-commercial educational use); Appendices B--C (tensors, index notation).
  \item M. E. Gurtin, E. Fried, L. Anand, \emph{The Mechanics and Thermodynamics of Continua}, Cambridge University Press, 2010. doi:10.1017/CBO9780511762956 (kinematics; Jacobi's formula).
  \item G. A. Holzapfel, \emph{Nonlinear Solid Mechanics: A Continuum Approach for Engineering}, Wiley, Chichester, 2000, Ch.~2 (kinematics).
\end{itemize}
\medskip
\textbf{Further reading.} P. A. Kelly, \emph{Solid Mechanics Part III: Foundations of Continuum Solid Mechanics}, University of Auckland lecture notes (online), kinematics chapter.
\end{frame}

  \end{document}